\chapter{Clustering}
\label{ch:clustering}

\chapterguide%
  {\chapref{ch:unsupervised}, feature scaling, vectors, distances, and the Gaussian distribution from \chapref{ch:foundations}.}
  {compare centroid, hierarchical, density, and soft clustering; assess stability and usefulness.}

This chapter applies the unsupervised workflow to \term{clustering}: finding groups, hierarchies, dense regions, or partial memberships.

Each method defines a group differently: $K$-means uses compact center-based groups, hierarchical methods nested structure, DBSCAN density connectivity, and mixtures probabilistic components.

Combine scores, stability, plots, and domain knowledge; an internal score alone does not make a clustering meaningful.

\begin{mltable}
\begin{tabularx}{\linewidth}{@{}>{\bfseries\leavevmode\color{MLNavy}}p{0.19\linewidth}p{0.34\linewidth}X@{}}
\toprule
\rowcolor{MLPaleBlue!92}
Assignment style & Meaning & Examples \\
\midrule
Hard clustering & Every row receives one cluster label. The memberships are mutually exclusive in the final partition. & $K$-means, $K$-medoids, DBSCAN, agglomerative clustering \\
Soft clustering & Every row can have partial membership in several groups, usually summarized by a vector of membership degrees or component responsibilities. & Fuzzy $C$-means, Gaussian mixture models \\
\bottomrule
\end{tabularx}
\end{mltable}

Hard and soft are about the \emph{form of the assignment}, not about whether an algorithm is better. If customers can genuinely sit between segments, a membership vector may preserve useful ambiguity that a single label throws away.

\section{$K$-means}
\label{sec:kmeans}

$K$-means divides observations into $K$ groups by minimizing squared distances from group centers:
\[
  \argmin_{S_1,\ldots,S_K}
  \sum_{j=1}^{K}\sum_{\bm{x}_i\in S_j}
  \|\bm{x}_i-\bm{\mu}_j\|_2^2 .
\]

In plain language, the formula says: choose $K$ centers, assign every row to a center, and move the centers until the rows are as close as possible to the center of their own group.

\begin{mlfigure}
\begin{tikzpicture}[scale=0.85, font=\small]
\foreach \p in {(0.4,0.6),(0.9,1.1),(0.2,1.3),(1.1,0.5),(0.7,1.6),(1.4,1.2)}
  \fill \p circle (2.2pt);
\node[cross out, draw, thick, minimum size=6pt] at (0.78,1.05) {};
\node at (0.78,2.1) {cluster 1};
\foreach \p in {(4.1,0.4),(4.7,0.9),(5.2,0.3),(4.5,1.4),(5.4,1.0),(4.0,1.1)}
  \fill[gray!60] \p circle (2.2pt);
\node[cross out, draw, thick, minimum size=6pt] at (4.65,0.85) {};
\node at (4.65,2.1) {cluster 2};
\foreach \p in {(2.2,3.3),(2.8,3.9),(3.3,3.2),(2.5,4.3),(3.5,3.9),(2.0,3.9)}
  \fill[gray!35] \p circle (2.2pt);
\node[cross out, draw, thick, minimum size=6pt] at (2.72,3.75) {};
\node at (4.7,3.75) {cluster 3};
\node[align=left] at (7.6,2.2) {$\times$ = cluster center $\bm{\mu}_j$\\(the mean of its points)};
\end{tikzpicture}
\end{mlfigure}

The algorithm (Lloyd's algorithm) alternates between two steps until nothing changes:

\begin{enumerate}
  \item \term{Assign}: send each observation to its nearest center.
  \item \term{Update}: move each center to the mean of the observations assigned to it.
\end{enumerate}

% === VISUAL ADDITION: kmeans-four-frame ===
\subsection{$K$-means as a four-frame loop}

Initialize centers, assign points, update centers, and repeat.

\begin{mlfigure}
\begin{tikzpicture}[font=\scriptsize,scale=0.78]
\def\pts{
  (0.45,0.60),(0.75,1.05),(1.15,0.75),(0.9,1.45),(1.45,1.25),
  (2.75,0.70),(3.15,1.00),(3.45,0.62),(3.05,1.55),(3.62,1.35)
}

% panel 1
\begin{scope}
  \draw[MLMuted!60] (0,0) rectangle (4.05,2.25);
  \foreach \p in \pts \fill[MLMuted] \p circle (2.1pt);
  \node[cross out,draw=MLBlue,very thick,minimum size=7pt] at (1.9,0.35) {};
  \node[cross out,draw=MLOrange,very thick,minimum size=7pt] at (2.15,1.95) {};
  \node[font=\scriptsize\bfseries\color{MLNavy}] at (2.03,2.58) {1. Initialize};
  \node[align=center,text=MLMuted] at (2.03,-0.35) {start with two centers};
\end{scope}

% panel 2
\begin{scope}[xshift=4.75cm]
  \draw[MLMuted!60] (0,0) rectangle (4.05,2.25);
  \foreach \p in {(0.45,0.60),(0.75,1.05),(1.15,0.75),(0.9,1.45),(1.45,1.25)}
    \fill[MLBlue] \p circle (2.1pt);
  \foreach \p in {(2.75,0.70),(3.15,1.00),(3.45,0.62),(3.05,1.55),(3.62,1.35)}
    \fill[MLOrange] \p circle (2.1pt);
  \node[cross out,draw=MLBlue,very thick,minimum size=7pt] at (1.9,0.35) {};
  \node[cross out,draw=MLOrange,very thick,minimum size=7pt] at (2.15,1.95) {};
  \node[font=\scriptsize\bfseries\color{MLNavy}] at (2.03,2.58) {2. Assign};
  \node[align=center,text=MLMuted] at (2.03,-0.35) {nearest center wins};
\end{scope}

% panel 3
\begin{scope}[xshift=9.50cm]
  \draw[MLMuted!60] (0,0) rectangle (4.05,2.25);
  \foreach \p in {(0.45,0.60),(0.75,1.05),(1.15,0.75),(0.9,1.45),(1.45,1.25)}
    \fill[MLBlue] \p circle (2.1pt);
  \foreach \p in {(2.75,0.70),(3.15,1.00),(3.45,0.62),(3.05,1.55),(3.62,1.35)}
    \fill[MLOrange] \p circle (2.1pt);
  \draw[-{Stealth[length=1.7mm]},MLBlue,thick] (1.9,0.35) -- (0.95,1.02);
  \draw[-{Stealth[length=1.7mm]},MLOrange,thick] (2.15,1.95) -- (3.20,1.04);
  \node[cross out,draw=MLBlue,very thick,minimum size=7pt] at (0.95,1.02) {};
  \node[cross out,draw=MLOrange,very thick,minimum size=7pt] at (3.20,1.04) {};
  \node[font=\scriptsize\bfseries\color{MLNavy}] at (2.03,2.58) {3. Update};
  \node[align=center,text=MLMuted] at (2.03,-0.35) {centers move to means};
\end{scope}

% panel 4, second row centered
\begin{scope}[xshift=4.75cm,yshift=-3.45cm]
  \draw[MLMuted!60] (0,0) rectangle (4.05,2.25);
  \foreach \p in {(0.45,0.60),(0.75,1.05),(1.15,0.75),(0.9,1.45),(1.45,1.25)}
    \fill[MLBlue] \p circle (2.1pt);
  \foreach \p in {(2.75,0.70),(3.15,1.00),(3.45,0.62),(3.05,1.55),(3.62,1.35)}
    \fill[MLOrange] \p circle (2.1pt);
  \node[cross out,draw=MLBlue,very thick,minimum size=7pt] at (0.95,1.02) {};
  \node[cross out,draw=MLOrange,very thick,minimum size=7pt] at (3.20,1.04) {};
  \node[font=\scriptsize\bfseries\color{MLNavy}] at (2.03,2.58) {4. Repeat};
  \node[align=center,text=MLGreen] at (2.03,-0.35) {stop when nothing changes};
\end{scope}

\draw[-{Stealth[length=1.8mm]},MLTeal,thick] (4.15,1.15) -- (4.62,1.15);
\draw[-{Stealth[length=1.8mm]},MLTeal,thick] (8.90,1.15) -- (9.37,1.15);
\draw[-{Stealth[length=1.8mm]},MLTeal,thick] (11.52,-0.35) .. controls (11.2,-1.6) and (8.2,-1.6) .. (7.0,-1.2);
\end{tikzpicture}
\end{mlfigure}
% === END VISUAL ADDITION ===

\begin{workedexample}
Cluster the one-dimensional points $\{1,2,3,10,11,12\}$ with $K=2$, starting from centers $c_1=1$ and $c_2=12$.
\emph{Assign:} points $1,2,3$ are nearer $c_1$; points $10,11,12$ are nearer $c_2$.
\emph{Update:} $c_1=(1+2+3)/3=2$ and $c_2=(10+11+12)/3=11$.
\emph{Assign again:} nothing moves (e.g.\ point 3 is at distance 1 from $c_1$ and 8 from $c_2$), so the algorithm stops. The objective is $(1+0+1)+(1+0+1)=4$, and the two obvious groups have been found in one round.
\end{workedexample}

Each round cannot increase the objective, but convergence may be local. Use $k$-means++ and multiple restarts (Lab~4 sets \texttt{n\_init} to at least 20). Mini-batches trade some fit quality for speed.

$K$-means implicitly assumes clusters are compact, roughly spherical, and similar in size; it struggles with elongated or irregular shapes, very different cluster sizes, and outliers (a single far-away point drags its center toward it).

\section{Choosing $K$: the elbow and the silhouette}

Use inertia, silhouette, stability, and domain evidence to choose $K$.

\begin{mlfigure}
\begin{tikzpicture}[baseline=(current bounding box.north)]
\begin{axis}[
  width=0.46\linewidth, height=4.3cm, axis lines=left,
  title style={font=\small\bfseries\color{MLNavy}}, title={Elbow: within-cluster distance},
  xlabel={$K$}, ylabel={Inertia},
  label style={font=\scriptsize\color{MLMuted}}, tick label style={font=\scriptsize\color{MLMuted}},
  xmin=0.5, xmax=9.5, ymin=0
]
\addplot[MLBlue,very thick,mark=*,mark size=1.5pt] coordinates
  {(1,980)(2,520)(3,290)(4,205)(5,178)(6,160)(7,146)(8,135)(9,126)};
\draw[MLRed,thick,dashed] (axis cs:4,0) -- (axis cs:4,980);
\node[font=\scriptsize\color{MLRed},anchor=west] at (axis cs:4.3,760) {the bend};
\end{axis}
\end{tikzpicture}
\hfill
\begin{tikzpicture}[baseline=(current bounding box.north)]
\begin{axis}[
  width=0.46\linewidth, height=4.3cm, axis lines=left,
  title style={font=\small\bfseries\color{MLNavy}}, title={Silhouette: separation quality},
  xlabel={$K$}, ylabel={Score},
  label style={font=\scriptsize\color{MLMuted}}, tick label style={font=\scriptsize\color{MLMuted}},
  xmin=0.5, xmax=9.5, ymin=0, ymax=0.8
]
\addplot[MLTeal,very thick,mark=square*,mark size=1.5pt] coordinates
  {(2,0.42)(3,0.51)(4,0.63)(5,0.55)(6,0.48)(7,0.41)(8,0.36)(9,0.31)};
\draw[MLRed,thick,dashed] (axis cs:4,0) -- (axis cs:4,0.8);
\node[font=\scriptsize\color{MLRed},anchor=west] at (axis cs:4.3,0.72) {the peak};
\end{axis}
\end{tikzpicture}
\end{mlfigure}

\begin{intuition}
Inertia decreases as $K$ increases, reaching zero with one cluster per point. Look for diminishing gains rather than its minimum. Silhouette favors compact, separated groups; neither diagnostic alone establishes usefulness.
\end{intuition}

\section{Initialization and stability}

Different starting centers can end in different solutions, with different inertia and different assignments.

Compare restarts and resamples. Very different solutions at the same $K$ indicate ambiguity worth reporting.

\begin{keyidea}
Check multiple initializations and resamples before interpreting clusters.
\end{keyidea}

\subsection{$K$-medoids}

$K$-medoids is a robust alternative to $K$-means: each cluster is represented by an actual observation (its \term{medoid}) rather than a mean, and any distance function may be used. Medoids resist outliers (a far point cannot drag a medoid the way it drags a mean) and make sense for data where averaging is meaningless --- at a higher computational price.

\section{Gaussian mixture models}

A Gaussian mixture model (GMM) is a probabilistic relative of $K$-means. It assumes that the observations were generated from $K$ Gaussian components with mixing weights $\pi_k$:
\[
  p(\bm{x})
  =\sum_{k=1}^{K}
  \pi_k\,
  \mathcal{N}(\bm{x}\given\bm{\mu}_k,\Sigma_k),
  \qquad
  \sum_{k=1}^{K}\pi_k=1 .
\]
Instead of hard assignments, each observation receives a \term{responsibility}: the posterior probability that component $k$ generated it,
\[
  \gamma_{ik}
  =\frac{\pi_k\,\mathcal{N}(\bm{x}_i\given\bm{\mu}_k,\Sigma_k)}
        {\sum_{j=1}^{K}\pi_j\,\mathcal{N}(\bm{x}_i\given\bm{\mu}_j,\Sigma_j)} .
\]
\term{Expectation-maximization (EM)} alternates responsibility estimation (E-step) with weighted updates of $\pi_k$, $\bm\mu_k$, and $\Sigma_k$ (M-step), improving data log-likelihood until convergence.

\begin{intuition}
EM alternates estimating memberships and component shapes. $K$-means follows a similar cycle with hard assignments.
\end{intuition}

GMMs provide soft memberships, elliptical components, and density-based anomaly scores (\chapref{ch:anomaly}). Responsibilities are probabilities under the mixture assumptions, not proof of real-world group membership.

\section{Fuzzy $C$-means: soft distance-based clustering}

\term{Fuzzy $C$-means} (FCM) keeps the centroid idea of $K$-means but replaces a one-hot assignment with membership degrees $u_{ik}\in[0,1]$ that sum to 1 across clusters for each observation. A common objective is
\[
  \sum_i\sum_{k=1}^{K} u_{ik}^{m}\,\|\bm{x}_i-\bm{c}_k\|_2^2,
\]
where $m>1$ is the \term{fuzziness parameter}. Larger $m$ produces softer memberships; $m=2$ is a common teaching default.

The algorithm alternates between two updates:
\begin{enumerate}
  \item recompute each centroid as a membership-weighted mean;
  \item recompute each observation's memberships from its relative distances to all centroids.
\end{enumerate}
The iterations stop when memberships or centers change by less than a chosen tolerance.

\begin{intuition}
$K$-means asks, ``Which center is closest?'' FCM asks, ``How strongly does each center describe this point?'' A boundary case might be 0.52 in cluster A and 0.45 in cluster B instead of being forced to pretend it is unambiguously A.
\end{intuition}

\begin{pitfall}
FCM memberships are distance-based degrees controlled by scale, cluster count, and fuzziness. Unlike GMM responsibilities, they do not imply a generative probability model.
\end{pitfall}

\section{Hierarchical clustering}
\label{sec:hierarchical-clustering}

\term{Agglomerative clustering} starts with one cluster per observation and repeatedly merges the closest pair. A \term{dendrogram} records the hierarchy.

\begin{mltable}
\begin{tabularx}{\linewidth}{@{}>{\bfseries\leavevmode\color{MLNavy}}p{0.22\linewidth}p{0.27\linewidth}X@{}}
\toprule
\rowcolor{MLPaleBlue!92}
Direction & Start & Repeated operation \\
\midrule
Agglomerative (bottom-up) & One cluster per observation & Merge the closest clusters until one hierarchy remains \\
Divisive (top-down) & One cluster containing all observations & Split a current cluster into smaller groups until a stopping rule is reached \\
\bottomrule
\end{tabularx}
\end{mltable}

Divisive clustering instead recursively splits clusters, sometimes starting from dissimilar seed observations.

\begin{workedexample}
For $\{1,2,4,8\}$, single linkage merges 1 and 2 at distance 1, adds 4 at distance 2, then adds 8 at distance 4:
\end{workedexample}

\begin{mlfigure}
\begin{tikzpicture}[font=\small, y=0.66cm, x=1.15cm]
\draw[-{Stealth[length=2mm]},MLMuted,thick] (-0.7,0) -- (-0.7,4.6);
\node[rotate=90,font=\scriptsize\color{MLMuted}] at (-1.55,2.3)
  {Merge distance};
\foreach \y in {1,2,4} {
  \draw[MLMuted] (-0.78,\y) -- (-0.62,\y);
  \node[anchor=east,font=\scriptsize\color{MLMuted}] at (-0.90,\y) {\y};
}
\node[below=2pt,font=\scriptsize] at (0,0) {1};
\node[below=2pt,font=\scriptsize] at (1,0) {2};
\node[below=2pt,font=\scriptsize] at (2,0) {4};
\node[below=2pt,font=\scriptsize] at (3,0) {8};
\draw[thick] (0,0) -- (0,1) -- (1,1) -- (1,0);
\draw[thick] (0.5,1) -- (0.5,2) -- (2,2) -- (2,0);
\draw[thick] (1.25,2) -- (1.25,4) -- (3,4) -- (3,0);
\draw[MLRed,thick,dashed] (-0.4,3) -- (3.30,3);
\node[anchor=west,align=left,font=\scriptsize\color{MLRed},
      fill=white,rounded corners=0.5mm,inner sep=2pt]
  at (3.42,3) {cut here\\$\rightarrow$ 2 clusters};
\end{tikzpicture}
\end{mlfigure}

Cut the dendrogram at a chosen height to obtain clusters. The \term{linkage} rule defines distance between clusters:

\begin{itemize}
  \item \term{single linkage} uses the closest pair of observations --- follows chains, finds elongated clusters, sensitive to noise bridges;
  \item \term{complete linkage} uses the farthest pair --- compact, similar-diameter clusters;
  \item \term{average linkage} uses the average pairwise distance --- a middle ground;
  \item \term{Ward linkage} merges whichever pair least increases within-cluster variance --- behaves most like $K$-means and is a common default.
\end{itemize}

Pairwise distances can be costly for large datasets; standard hierarchical clustering has no built-in new-point assignment.

\section{DBSCAN: clustering by density}
\label{sec:dbscan}

DBSCAN defines clusters as regions where points are packed densely, using two hyperparameters: a neighborhood radius $\varepsilon$ and a minimum count \term{minPts}. Every point gets one of three roles:

\begin{description}[leftmargin=1.05in,style=nextline]
  \item[Core point] Has at least minPts points within distance $\varepsilon$ --- it sits in a dense region.
  \item[Border point] Is within $\varepsilon$ of a core point but is not dense enough itself --- the edge of a cluster.
  \item[Noise point] Neither --- it belongs to no cluster at all.
\end{description}

\begin{mlfigure}
\begin{tikzpicture}[scale=1.5, font=\small]
\foreach \p in {(1.0,1.0),(1.35,1.15),(0.8,1.4),(1.2,0.7),(1.55,0.8),(0.65,0.85),(1.05,1.45),(1.5,1.35)}
  \fill \p circle (1.6pt);
\fill (1.9,1.5) circle (1.6pt);
\fill (3.4,0.7) circle (1.6pt);
\draw (1.0,1.0) circle (0.55);
\draw[-{Stealth[length=1.6mm]}] (1.0,1.0) -- (0.611,0.611)
  node[below left=-2pt] {$\varepsilon$};
\draw[dashed] (1.9,1.5) circle (0.55);
\draw[dashed] (3.4,0.7) circle (0.55);
\node[align=center, below] at (1.0,0.32) {core point:\\$\geq$ minPts neighbors};
\node[align=center, above] at (2.15,2.08) {border point: sparse itself,\\but near a core point};
\node[align=center, below] at (3.4,0.02) {noise point};
\end{tikzpicture}
\end{mlfigure}

Connected core points form clusters, border points attach, and noise remains unassigned. DBSCAN handles irregular shapes but struggles with varying densities or high dimensions. Use a sorted neighbor-distance plot to guide $\varepsilon$. In Lab~4, changing $\varepsilon$ from 0.75 to 0.5 changes one cluster with 14 noise points into five clusters with roughly one-third noise.

\section{Which clustering method?}

\begin{mltable}
\begin{tabular}{@{}p{0.16\linewidth}p{0.12\linewidth}p{0.21\linewidth}p{0.14\linewidth}p{0.22\linewidth}@{}}
\toprule
\rowcolor{MLPaleBlue!92}
Method & Needs $K$? & Cluster shapes & Outliers & Notable trait \\
\midrule
$K$-means & yes & compact, spherical & distort means & fast, scalable, simple \\
GMM & yes & elliptical & distort fits & soft memberships, density \\
Hierarchical & no (cut later) & depends on linkage & can distort merges & full dendrogram of groupings \\
DBSCAN & no & arbitrary & labeled as noise & density-based, finds outliers \\
\bottomrule
\end{tabular}
\end{mltable}

$K$-means and GMM also assign new observations naturally (nearest center; highest responsibility); hierarchical clustering and DBSCAN have no built-in rule for unseen points, so a small classifier trained on the cluster labels often fills that gap in practice.

\section{Evaluating a clustering}

The \term{silhouette coefficient} for observation $i$ compares cohesion with separation:
\[
  s(i)=\frac{b(i)-a(i)}{\max\{a(i),b(i)\}},
\]
$a(i)$ is mean within-cluster distance; $b(i)$ is the smallest mean distance to another cluster. Near 1 means separation, near 0 a boundary, and negative values possible misassignment. Lab~4 reports the mean silhouette.

When reference labels happen to exist (for validation only), agreement measures such as the \term{adjusted Rand index} compare the discovered clusters with the labeled groups while correcting for chance.

\begin{keyidea}
Treat clusters as hypotheses. Support interpretation with scores, stability, and domain evidence.
\end{keyidea}

\begin{modelcard}{Clustering - Quick Guide}
\begin{tabularx}{\linewidth}{@{}>{\bfseries\leavevmode\color{MLNavy}}p{0.2\linewidth}X X@{}}
\toprule
\rowcolor{MLPaleBlue!92}
Method & Good for & Main caution \cr
\midrule
$K$-means & Compact, roughly spherical clusters; large numeric datasets & Requires $K$, scaling, and sensitivity to initialization/outliers \cr
Gaussian mixture & Overlapping elliptical clusters and soft membership & Distribution assumptions and local optima \cr
Hierarchical & Dendrograms and nested group structure & Distance/linkage choices and cost on large data \cr
DBSCAN & Irregular shapes and explicit noise points & Density parameters and varying-density clusters \cr
\bottomrule
\end{tabularx}
\end{modelcard}

\begin{quickcheck}
\begin{enumerate}
  \item What kind of cluster geometry does \(K\)-means favor, and why can scaling change the resulting partition?
  \item What does DBSCAN do with observations that do not belong to a sufficiently dense region?
\end{enumerate}
\end{quickcheck}
